\chapter{Observabilité et réponse aux incidents}\label{ch:pl:observability-and-incident-response}

Le 22 août 2013, le système qui diffuse les cotations consolidées des actions cotées au Nasdaq commence à défaillir à 11~h~48 du matin, avec
des interruptions intermittentes sur certains de ses canaux. À 12~h~03, il cesse complètement de diffuser des cotations~; une suspension réglementaire arrête la négociation de ces
actions sur toutes les plateformes~; la panne technique est réparée à 12~h~24, et la cotation comme la négociation reprennent à 15~h~25. La supervision
d'une société de trading ce matin-là pouvait sembler parfaite~: chaque processus était actif, chaque file vide, chaque heartbeat à l'heure~--- parce que rien n'arrivait. Ce chapitre construit la
supervision qui aurait détecté le problème~: des signaux portant sur ce que reçoivent les utilisateurs de la plateforme plutôt que sur les processus qui les servent,
des objectifs assortis de \omterm{def:pl:observability-and-incident-response:budget}{budgets d'erreur}, des alertes qui ne déclenchent un appel que lorsqu'une personne doit agir, et la pratique d'incident qui suit l'appel.

\section{Métriques, journaux et traces}

\begin{definition}[Observabilité]\label{def:pl:observability-and-incident-response:observability}
L'\emph{observabilité}\index{observabilité} est le degré auquel l'état interne d'un système en fonctionnement peut être déduit de
ce qu'il émet~--- ses \omterm{def:pl:observability-and-incident-response:metric}{métriques}, ses journaux et ses traces~--- sans déployer de nouveau code pour poser une nouvelle question.
\end{definition}

\begin{definition}[Métrique, journalisation structurée]\label{def:pl:observability-and-incident-response:metric}
Une \emph{métrique}\index{métrique} est une mesure numérique nommée, munie d'étiquettes, échantillonnée ou agrégée dans le temps~: un compteur (messages reçus), une
jauge (profondeur de file), un histogramme (latence, comme dans le livre~13, chapitre~5). La \emph{journalisation structurée}\index{journalisation structurée} écrit chaque événement de journal
sous la forme d'un enregistrement de champs nommés~--- heure, niveau, service, événement, identifiants, valeurs~--- plutôt que sous la forme d'une phrase, de sorte que les journaux peuvent être interrogés
comme des données.
\end{definition}

\begin{definition}[Traçage distribué, span]\label{def:pl:observability-and-incident-response:trace}
Le \emph{traçage distribué}\index{traçage distribué} suit une unité de travail~--- une requête, un message de données de marché, un ordre~--- à travers
les services qu'elle traverse. Chaque étape est un \emph{span}\index{span}, une unité de travail dotée d'un début, d'une fin, d'un identifiant de trace partagé par toutes les
étapes et de l'identifiant de l'étape qui l'a causée, de sorte que le chemin peut être reconstruit sous forme d'arbre.
\end{definition}

La chaîne du chapitre est celle de la plateforme elle-même~: la capture (chapitre~\ref{ch:pl:capturing-and-storing-tick-data}), la \omterm{def:pl:a-tick-store-on-open-formats:store}{base de ticks}, le \omterm{def:pl:position-and-pnl-services:service}{service de positions}
(chapitre~\ref{ch:pl:position-and-pnl-services}) et le risque en temps réel (chapitre~\ref{ch:pl:real-time-risk}). \texttt{firm.observe} fournit
à chacun un registre de \omterm{def:pl:observability-and-incident-response:metric}{métriques}, des enregistrements de journal structurés et des \omterm{def:pl:observability-and-incident-response:trace}{spans} au sens d'OpenTelemetry, transmis d'un service au suivant. Un message tracé
montre où passe son temps (\cref{fig:pl:observability-and-incident-response:trace})~: une fraction de milliseconde dans la capture, le stockage et
les positions, puis l'attente dans la file du service de risque~--- 50~millisecondes un jour normal, une minute lorsque le service de risque prend du retard.

\omcode{code/firm/observe/firm_observe.py}{103}{123}{Spans~: un span racine démarre une trace, chaque enfant porte l'identifiant de la trace et celui de son
parent, et une trace est reconstruite à partir d'eux dans l'ordre chronologique.}\label{lst:pl:observability-and-incident-response:tracer}

\begin{omfigure}
\begin{tikzpicture}
\begin{semilogxaxis}[width=11cm,height=5.2cm,xbar,bar width=6pt,xmin=0.005,xmax=200000,y dir=reverse,ytick={0,1,2,3,4},
  yticklabels={capture,stockage,positions,file du risque,calcul de risque},xlabel={durée du span (ms, échelle logarithmique)},grid=major,grid style={gray!20},
  tick label style={font=\small},label style={font=\small},table/col sep=comma,log origin=infty,
  legend style={font=\small,at={(0.5,1.03)},anchor=south,legend columns=2,draw=none,
    /tikz/every even column/.append style={column sep=4mm}}]
\addplot[fill=omDef!40,draw=omDef] table[y=row,x=dur_ms]{figdata/platforms/28-observability-and-incident-response/trace_normal.csv};
\addplot[fill=omThm!50,draw=omThm] table[y=row,x=dur_ms]{figdata/platforms/28-observability-and-incident-response/trace_slow.csv};
\legend{journée normale,consommateur de risque lent}
\end{semilogxaxis}
\end{tikzpicture}

\omcaption{Les \omterm{def:pl:observability-and-incident-response:trace}{spans} d'un message de données de marché tracé à travers la chaîne (échelle logarithmique). La capture, le stockage et les positions prennent 0.02, 0.2 et
0.1~ms~; le service de risque calcule en 1~ms. Ce qui change lorsque le service de risque prend du retard, c'est sa file~: 50~ms un jour normal, 60~secondes
au plus fort de l'incident du consommateur lent. Données~: \texttt{fig\_observe.py}.}\label{fig:pl:observability-and-incident-response:trace}
\end{omfigure}

Chacun des trois répond à une question différente. Les \omterm{def:pl:observability-and-incident-response:metric}{métriques} disent que quelque chose ne va pas, à faible coût, pour tout, en permanence. Les journaux disent ce qui
s'est passé dans un service. Les traces disent où le temps est passé à travers les services, pour les unités de travail échantillonnées. La plupart des sociétés de trading disposent des deux premiers
et pas du troisième, et n'apprennent où un message a passé sa minute qu'en lisant les journaux de quatre services côte à côte, avec quatre horloges.

\section{Niveaux de service et budgets d'erreur}

\begin{definition}[Indicateur de niveau de service, objectif de niveau de service]\label{def:pl:observability-and-incident-response:sli}
Un \emph{indicateur de niveau de service}\index{indicateur de niveau de service} (SLI) est une mesure soigneusement définie d'un aspect du service que reçoit un
utilisateur~: ici, la part des secondes pendant lesquelles les données de marché utilisées par le service de risque ont au plus cinq secondes. Un \emph{objectif de
niveau de service}\index{objectif de niveau de service} (SLO) est une cible fixée pour un SLI sur une période~: ici, 99.9~\% des secondes de négociation d'un mois de 30 jours.
\end{definition}

Le choix de l'indicateur se fait du côté de l'utilisateur. Un heartbeat du service de risque indique que le processus est vivant~; l'âge de ses données indique
s'il fait son travail. Le matin de l'incident, chaque heartbeat était à l'heure et chaque âge augmentait d'une seconde par seconde. Un SLO se distingue
de l'accord de niveau de service (SLA) du livre~14 par la partie qu'il engage~: un accord est une promesse faite à un client, assortie de pénalités, alors qu'un objectif est la cible propre de la société,
fixée plus strictement que toute promesse qu'elle fait.

\begin{definition}[Budget d'erreur]\label{def:pl:observability-and-incident-response:budget}
Le \emph{budget d'erreur}\index{budget d'erreur} d'un SLO est la quantité de mauvais service qu'il autorise sur sa période~--- 1~moins la cible, multiplié par les
événements~--- que le service consomme pour les incidents, la maintenance et les changements risqués~; lorsqu'il est épuisé, le travail de fiabilité prime sur
les fonctionnalités.
\end{definition}

Un objectif de 99.9~\% sur 30 jours de négociation de 23\,400 secondes autorise 702 mauvaises secondes par mois. Les journées calmes du chapitre~--- la latence normale
de la chaîne, de quelques dizaines de millisecondes, plus des pauses (une rafale, une collecte, un disque lent) qui surviennent 0.8 fois par heure et pendant lesquelles l'âge augmente
d'une seconde par seconde~--- en consomment 115, soit 16~\% du budget, en 18 pauses de plus de cinq secondes
(\cref{fig:pl:observability-and-incident-response:quiet}).

\begin{omfigure}
\begin{tikzpicture}
\begin{semilogyaxis}[width=11.5cm,height=5.2cm,xlabel={heure de la journée},ylabel={âge maximal (s, log)},
  grid=major,grid style={gray!20},tick label style={font=\small},label style={font=\small},table/col sep=comma,
  xmin=0,xmax=390,ymin=0.03,ymax=200,xtick={0,90,180,270,360},xticklabels={09~h~30,11~h~00,12~h~30,14~h~00,15~h~30}]
\addplot[omDef,thin,const plot] table[x=minute,y=max_age]{figdata/platforms/28-observability-and-incident-response/quiet.csv};
\draw[omThm,dashed] (axis cs:0,5) -- (axis cs:390,5) node[pos=0.02,above right,font=\scriptsize]{5 s~: la limite du SLI};
\draw[gray,dashed] (axis cs:0,30) -- (axis cs:390,30) node[pos=0.98,below left,font=\scriptsize]{30 s};
\draw[gray,dashed] (axis cs:0,60) -- (axis cs:390,60) node[pos=0.98,above left,font=\scriptsize]{60 s};
\end{semilogyaxis}
\end{tikzpicture}

\omcaption{Une journée calme du modèle~: l'âge maximal des données de marché utilisées par le service de risque à chaque minute. Hors des pauses, l'âge maximal d'une minute
reste d'environ 0.1 à 0.3~seconde, soit quelques fois les 50~ms typiques de la chaîne~; quelques pauses d'une ou deux secondes passent inaperçues~; une pause à 14~h~35 dure 56 secondes, dépassant la limite de cinq secondes du SLI et tout seuil statique
inférieur à une minute. Données~: \texttt{fig\_observe.py}.}\label{fig:pl:observability-and-incident-response:quiet}
\end{omfigure}

\section{Alerter pour que les personnes puissent agir}

Un appel réveille une personne. Il doit signifier qu'une personne doit agir maintenant, et il doit arriver assez tôt pour que l'action compte. Des
seuils statiques sur l'âge des données opposent les deux exigences~; une règle fondée sur le \omterm{def:pl:observability-and-incident-response:budget}{budget d'erreur} n'a pas à le faire.

\begin{definition}[Alerte de vitesse de consommation]\label{def:pl:observability-and-incident-response:burn}
Une \emph{alerte de vitesse de consommation}\index{alerte de vitesse de consommation} (burn-rate alert) déclenche un appel lorsque le \omterm{def:pl:observability-and-incident-response:budget}{budget d'erreur} est consommé plus vite qu'un multiple donné du rythme qui
l'épuiserait exactement sur la période du SLO, mesurée sur une longue fenêtre et confirmée sur une courte~: pour un objectif de 99.9~\%, le
SRE workbook de Google recommande de déclencher un appel à 14.4 fois la vitesse sur une heure (confirmée sur cinq minutes) ou à 6 fois la vitesse sur six heures (sur trente minutes).
\end{definition}

\omcode{code/firm/observe/firm_observe.py}{159}{176}{Alerte multi-fenêtres de vitesse de consommation~: la part des secondes défaillantes sur les fenêtres longue et
courte de chaque règle, comparée au rythme multiplié par la part du budget d'erreur.}\label{lst:pl:observability-and-incident-response:burn}

Trois incidents sont rejoués un jour calme, aux heures de l'horloge du crochet~: un arrêt du flux à partir de 12~h~03, qui dure 21 minutes~; des arrêts intermittents de 20
secondes chaque minute à partir de 11~h~48, pendant un quart d'heure~; et un consommateur lent à partir de 14~h~00, un service de risque dont le retard croît de 0.05 seconde
à chaque seconde. Cinq règles les surveillent (la \cref{tab:pl:observability-and-incident-response:rules}).

\begin{table}[htbp]
\centering\small
\begin{tabular}{@{}>{\raggedright\arraybackslash}p{3.6cm}>{\centering\arraybackslash}p{2.6cm}>{\centering\arraybackslash}p{2.6cm}>{\centering\arraybackslash}p{2.6cm}>{\centering\arraybackslash}p{2.6cm}@{}}
\toprule
règle & arrêt du flux & arrêts intermittents & consommateur lent & fausses alertes par mois\\
\midrule
heartbeat du processus & jamais & jamais & jamais & 0\\
âge des données supérieur à 5~s & 5~s & 5~s & 100~s & 18\\
âge des données supérieur à 30~s & 30~s & jamais & 600~s & 1\\
âge des données supérieur à 60~s & 60~s & jamais & 1\,200~s & 0\\
vitesse de consommation, règles à 1~h et 6~h & 56~s & 191~s & 151~s & 1\\
\bottomrule
\end{tabular}
\caption{\omterm{def:pl:observability-and-incident-response:ttd}{Délai de détection} de chaque incident rejoué, et appels sur trente jours calmes, pour cinq règles d'alerte. Le heartbeat ne se déclenche
jamais~; les seuils statiques achètent la rapidité au prix de fausses alertes ou manquent les arrêts intermittents~; les règles de vitesse de consommation détectent
les trois incidents en environ trois minutes et ne déclenchent qu'un appel par mois les jours calmes, pour la pause de 56 secondes de la journée.}
\label{tab:pl:observability-and-incident-response:rules}
\end{table}

\begin{definition}[Délai de détection, délai de rétablissement]\label{def:pl:observability-and-incident-response:ttd}
Le \emph{délai de détection}\index{délai de détection} d'un incident est le temps écoulé entre son début et la première alerte ou le premier signalement sur lequel
une personne agit~; son \emph{délai de rétablissement}\index{délai de rétablissement} est le temps écoulé entre son début et le moment où le service que reçoivent
les utilisateurs est de nouveau dans son objectif.
\end{definition}

Le \omterm{def:pl:observability-and-incident-response:ttd}{délai de rétablissement} se mesure sur le service que reçoit un utilisateur, à partir du début de l'incident~; le temps moyen de réparation (MTTR) du livre~14 est une
propriété d'un composant comme un circuit, moyennée sur ses pannes, depuis la panne jusqu'à la réparation.

Les seuils statiques montrent le compromis (la \cref{tab:pl:observability-and-incident-response:rules}). À cinq secondes, ils détectent tout
d'un coup et déclenchent 18 appels par mois pour des pauses qui se résorbent d'elles-mêmes~; en quelques semaines, personne n'y répond. À trente ou soixante secondes, les
fausses alertes disparaissent, mais les arrêts intermittents --- la phase de 11~h~48 du crochet --- n'atteignent jamais le seuil, et un consommateur lent
met dix ou vingt minutes à être repéré. Les règles de vitesse de consommation déclenchent un appel 56 secondes après le début de l'arrêt, 3 minutes après celui
des arrêts intermittents et 2 minutes et demie après celui du consommateur lent (la \cref{fig:pl:observability-and-incident-response:slow}), avec un seul appel par mois les jours calmes~: la pause de 56 secondes,
qui consomme bien une part réelle du budget et mérite un examen.

\begin{omfigure}
\begin{tikzpicture}
\begin{axis}[width=11.5cm,height=5.4cm,xlabel={minutes après le début de l'incident},ylabel={âge des données (s)},
  grid=major,grid style={gray!20},tick label style={font=\small},label style={font=\small},table/col sep=comma,
  xmin=-5,xmax=25,ymin=0,ymax=80,axis y line*=left]
\addplot[omDef,thick] table[x=t_min,y=age]{figdata/platforms/28-observability-and-incident-response/incident_slow.csv};
\draw[black,dashed] (axis cs:100/60,0) -- (axis cs:100/60,80) node[pos=0.97,left=3pt,font=\scriptsize]{5~s~: 1.7~min};
\draw[omThm,thick] (axis cs:151/60,0) -- (axis cs:151/60,80) node[pos=0.85,right,font=\scriptsize]{vitesse de consommation~: 2.5~min};
\draw[black,dashed] (axis cs:10,0) -- (axis cs:10,80) node[pos=0.97,left,font=\scriptsize]{30~s~: 10~min};
\draw[black,dashed] (axis cs:20,0) -- (axis cs:20,80) node[pos=0.97,left,font=\scriptsize]{60~s~: 20~min};
\node[omDef,font=\scriptsize,anchor=south east] at (axis cs:24.5,4) {âge des données du service de risque};
\end{axis}
\end{tikzpicture}

\omcaption{L'incident du consommateur lent~: à partir de 14~h~00, le service de risque traite les messages moins vite qu'ils n'arrivent, et l'âge de ses
données croît de 0.05 seconde à chaque seconde. Son processus reste vivant tout du long. Le seuil de cinq secondes se déclenche à 1.7 minute, la règle de vitesse de consommation
à 2.5 minutes, les seuils de 30 et 60 secondes après 10 et 20 minutes. Données~: \texttt{fig\_observe.py}.}\label{fig:pl:observability-and-incident-response:slow}
\end{omfigure}

\begin{omfigure}
\begin{tikzpicture}
\begin{axis}[width=11.5cm,height=5.4cm,xlabel={minutes après 11~h~48},ylabel={consommation (fois le budget)},
  grid=major,grid style={gray!20},tick label style={font=\small},label style={font=\small},table/col sep=comma,
  xmin=-5,xmax=15,ymin=0,ymax=270]
\addplot[gray,thin] table[x=t_min,y=burn_5m]{figdata/platforms/28-observability-and-incident-response/incident_intermittent.csv};
\addplot[omThm,thick] table[x=t_min,y=burn_1h]{figdata/platforms/28-observability-and-incident-response/incident_intermittent.csv};
\draw[black,dashed] (axis cs:-5,14.4) -- (axis cs:15,14.4) node[pos=0.02,above right,font=\scriptsize]{14.4};
\draw[omThm,dotted,thick] (axis cs:191/60,0) -- (axis cs:191/60,270) node[pos=0.3,right,font=\scriptsize]{appel~: 3.2~min};
\node[gray,font=\scriptsize,anchor=north] at (axis cs:10,240) {fenêtre de 5 minutes};
\node[omThm,font=\scriptsize,anchor=south] at (axis cs:11,66) {fenêtre d'une heure};
\end{axis}
\end{tikzpicture}

\omcaption{Les arrêts intermittents~: 20 secondes sans données chaque minute à partir de 11~h~48. L'âge des données ne dépasse jamais 20 secondes, donc les seuils à 30
et 60 secondes ne se déclenchent jamais~; la vitesse de consommation sur la dernière heure dépasse 14.4 --- un cinquantième du budget du mois en une heure --- après 3.2 minutes,
la fenêtre de cinq minutes la confirmant. Données~: \texttt{fig\_observe.py}.}\label{fig:pl:observability-and-incident-response:intermittent}
\end{omfigure}

\section{Astreinte et commandement d'incident}

\begin{definition}[Rotation d'astreinte, procédure d'exploitation]\label{def:pl:observability-and-incident-response:oncall}
Une \emph{rotation d'astreinte}\index{rotation d'astreinte} désigne, selon un planning publié, les personnes qui répondent aux appels d'un service à chaque heure,
avec un suppléant derrière le titulaire. Une \emph{procédure d'exploitation}\index{procédure d'exploitation} (runbook) est la procédure écrite pour une alerte ou une panne connue~: ce qu'elle signifie,
comment la confirmer, quoi faire en premier, qui appeler.
\end{definition}

\begin{definition}[Responsable d'incident]\label{def:pl:observability-and-incident-response:ic}
Le \emph{responsable d'incident}\index{responsable d'incident} (incident commander) tient l'état global d'un incident et structure la réponse --- en désignant
qui modifie le système, qui communique, qui planifie --- sans effectuer lui-même le travail technique.
\end{definition}

Dans une société de trading, les premières minutes d'un incident concernent l'exposition, non le diagnostic~: la société opère-t-elle sur des données périmées, et faut-il arrêter
les stratégies~? La \omterm{def:pl:observability-and-incident-response:oncall}{procédure d'exploitation} d'une alerte d'âge des données commence donc par les kill switches du livre~13 et les limites de risque du livre~11, puis par
le diagnostic, et la première question du \omterm{def:pl:observability-and-incident-response:ic}{responsable d'incident} au desk de trading est de savoir ce que la société détient. Un seul groupe modifie le système pendant un
incident, afin que les correctifs n'entrent pas en collision~; une seule personne communique, afin que le desk, le risque et la conformité entendent la même chose.

\begin{table}[htbp]
\centering\small
\begin{tabular}{lll}
\toprule
heure & source & événement\\
\midrule
12:03:00 & log & tickcap~: dernier message sur les canaux de cotation\\
12:03:05 & alerte & âge des données supérieur à 5~s~: appel\\
12:03:30 & alerte & âge des données supérieur à 30~s~: appel\\
12:03:56 & alerte & vitesse de consommation~: appel\\
12:04:00 & alerte & âge des données supérieur à 60~s~: appel\\
12:24:00 & log & tickcap~: reprise des canaux de cotation\\
12:24:00 & log & rtrisk~: âge des données repassé sous 5~s\\
\bottomrule
\end{tabular}
\caption{Chronologie de l'arrêt du flux rejoué, fusionnée à partir des alertes et des journaux structurés par \texttt{firm.observe.timeline}. \omterm{def:pl:observability-and-incident-response:ttd}{Délai de détection}~:
56~secondes avec la règle de vitesse de consommation~; \omterm{def:pl:observability-and-incident-response:ttd}{délai de rétablissement}~: 21 minutes.}\label{tab:pl:observability-and-incident-response:timeline}
\end{table}

\section{Un incident rejoué, et la revue qui suit}

\begin{definition}[Revue sans recherche de coupable]\label{def:pl:observability-and-incident-response:blameless}
Une \emph{revue sans recherche de coupable}\index{revue sans recherche de coupable} (blameless review~; post-incident review, postmortem) reconstitue la chronologie d'un incident et ses causes contributives
sans désigner aucun individu ni aucune équipe, en partant du principe que chacun a agi de bonne foi avec les informations dont il disposait, et
se termine par des actions qui modifient le système plutôt que par des exhortations à la prudence.
\end{definition}

La revue de l'arrêt rejoué s'écrit presque d'elle-même à partir du \cref{tab:pl:observability-and-incident-response:timeline}. Les données se sont arrêtées à 12~h~03~; la
\omterm{def:pl:observability-and-incident-response:burn}{alerte de vitesse de consommation} est arrivée à 12:03:56~; le flux est revenu à 12~h~24, et les données du service de risque étaient de nouveau fraîches dans la même seconde. Les questions qu'elle pose ne
portent pas sur qui a failli, mais sur ce que le système a laissé se produire~: pourquoi les tableaux de bord de la société affichaient du vert (ils mesuraient les processus, non les données)~; ce que les stratégies
ont fait des prix périmés pendant la première minute (le déclencheur du kill switch portait sur la santé des processus, non sur l'âge des données)~; pourquoi la phase intermittente à partir de
11~h~48 est passée inaperçue (les seuils étaient supérieurs à ses vingt secondes). Ses actions sont des modifications~: un SLI sur l'âge des données pour chaque consommateur, le déclenchement des \omterm{def:pl:observability-and-incident-response:burn}{alertes de vitesse de consommation},
un kill switch sur les données périmées, et le rejeu lui-même conservé comme test (chapitre~\ref{ch:pl:testing-and-continuous-delivery}) afin que la prochaine
modification de l'alerte soit vérifiée contre lui. Un incident qui devient un événement de perte opérationnelle (livre~6, chapitre~28) entre aussi dans les données de
pertes de la société~; c'est la revue qui rend ces données utiles.

\begin{dated}{2026-09}{L'incident SIP de 2013 et les pratiques SRE}\label{dat:pl:observability-and-incident-response:sip}
L'alerte fournisseur de Nasdaq sur le problème du SIP UTP du 22 août 2013 donne la chronologie~: des pannes intermittentes sur plusieurs canaux de cotation à partir de
11~h~48, une panne sur tous les canaux de cotation à 12~h~03, le problème résolu et l'environnement stable à 12~h~24, la connectivité rétablie à partir de
14~h~00, et toute la cotation et le trading des titres cotés sur Nasdaq ont repris à 15~h~25~; les flux de transactions n'ont pas été touchés, et un arrêt réglementaire a été
décidé pour toutes les négociations de ces titres. \emph{Site Reliability Engineering} de Google définit les SLI, les SLO et le \omterm{def:pl:observability-and-incident-response:budget}{budget d'erreur}, et décrit
le commandement d'incident et les revues post-mortem sans recherche de coupable~; son workbook recommande, pour un objectif de 99.9~\%, de déclencher des alertes sur des vitesses de consommation de 14.4 sur une heure et de 6
sur six heures, chacune confirmée sur une fenêtre douze fois plus courte. OpenTelemetry définit un \omterm{def:pl:observability-and-incident-response:trace}{span} comme une unité de travail portant les identifiants de trace, de \omterm{def:pl:observability-and-incident-response:trace}{span}
et de parent.
\end{dated}

\section{Tutoriel~: trois heures de silence}\label{tut:pl:observability-and-incident-response:tut}

\begin{tutorial}
\textbf{Objectif.} Instrumenter la chaîne, définir son SLO, rejouer trois incidents, comparer les règles d'alerte et construire la chronologie de l'incident.
\textbf{État final~:} les \cref{tab:pl:observability-and-incident-response:rules,tab:pl:observability-and-incident-response:timeline}.
\begin{tutsteps}
\item \textbf{Traces}~: \texttt{pl\_observe.traced(ts, lag)} un jour normal et pendant le consommateur lent.
\item \textbf{Jours calmes}~: \texttt{quiet\_day(seed)} et \texttt{budget()}, le \omterm{def:pl:observability-and-incident-response:budget}{budget d'erreur} consommé dans un mois.
\item \textbf{Incidents}~: \texttt{incident(kind)} pour les trois types.
\item \textbf{Règles}~: \texttt{compare()}~: \omterm{def:pl:observability-and-incident-response:ttd}{délai de détection} et fausses alertes pour chaque règle.
\item \textbf{Chronologie}~: \texttt{stall\_timeline()} et la revue.
\end{tutsteps}
\textbf{Pour aller plus loin.} Ajouter une règle de ticket (vitesse de consommation de 1 sur trois jours) et compter ce qu'elle détecte et que les alertes ne détectent pas~; rendre les
pauses deux fois plus fréquentes et voir quelles règles survivent.
\end{tutorial}

\section{Construction~: observabilité}\label{bld:pl:observability-and-incident-response:observe}

\begin{build}
\bfield{Rôle} Savoir ce que reçoivent les utilisateurs de la plateforme, n'appeler une personne que lorsqu'une action est nécessaire, et reconstituer tout incident à partir de
ce que la plateforme a émis.
\bfield{Interface} \texttt{Registry} (\texttt{counter}, \texttt{gauge}, \texttt{histogram}), \texttt{log}, \texttt{Tracer} (\texttt{start},
\texttt{finish}, \texttt{trace}), \texttt{SLO}, \texttt{error\_budget\_used}, \texttt{threshold\_pages}, \texttt{burn\_rate\_pages},
\texttt{timeline}.
\bfield{Règles} SLI mesurés du côté de l'utilisateur~; chaque service émet des \omterm{def:pl:observability-and-incident-response:metric}{métriques}, des journaux structurés et des \omterm{def:pl:observability-and-incident-response:trace}{spans} avec des identifiants propagés~; des alertes
déclenchées par les vitesses de consommation des SLO, des tickets pour les consommations lentes~; chaque alerte a une \omterm{def:pl:observability-and-incident-response:oncall}{procédure d'exploitation} (runbook)~; chaque incident a un \omterm{def:pl:observability-and-incident-response:ic}{responsable d'incident} et une \omterm{def:pl:observability-and-incident-response:blameless}{revue sans recherche de coupable} accompagnée
de modifications.
\bfield{Tests d'acceptation} \texttt{code/firm/observe/tests/}~: types et étiquettes du registre, une ligne de journal structuré~; une trace reconstruite à partir des liens parent
liens~; budgets et vitesses de consommation~; alertes à seuil avec un écart minimal~; une \omterm{def:pl:observability-and-incident-response:burn}{alerte de vitesse de consommation} à la 52\up{e} seconde défaillante d'une panne, et aucune avant~; une
chronologie fusionnée.
\bfield{Pour aller plus loin} Exemplaires reliant un bucket d'histogramme à une trace~; échantillonnage de traces fondé sur la queue~; alertes sur le CUSUM d'un SLI (livre~7,
chapitre~13).
\end{build}

\begin{omsources}
\item Nasdaq, \emph{UTP Vendor Alert \#2013-9: UTP SIP Issue on Thursday, August 22, 2013}.
\item B.~Beyer, C.~Jones, J.~Petoff et N.~R.~Murphy (dir.), \emph{Site Reliability Engineering}, O'Reilly, 2016 (chapitres 3, 4, 14, 15);
\emph{The Site Reliability Workbook}, 2018 (\emph{Alerting on SLOs}).
\item Documentation OpenTelemetry, \emph{Traces}.
\item One Quant Book~13, chapitres~5 et~24 (histogrammes de latence, kill switches)~; livre~14, chapitre~15 (accords de niveau de service (SLA), temps moyen de
réparation).
\end{omsources}

\section{Exercices}

\begin{exercise}[$\star$]\label{exo:pl:observability-and-incident-response:1}
Combien de secondes défaillantes un objectif de 99.9~\% autorise-t-il sur trente jours de 23\,400 secondes, et combien le mois calme en a-t-il consommé~?
\end{exercise}

\begin{exercise}[$\star$]\label{exo:pl:observability-and-incident-response:2}
Pourquoi le heartbeat du processus ne se déclenche jamais dans aucun des trois incidents~?
\end{exercise}

\begin{exercise}[$\star$]\label{exo:pl:observability-and-incident-response:3}
Combien de temps la règle de vitesse de consommation met-elle à déclencher une alerte pour un arrêt complet, et d'où vient ce nombre~?
\end{exercise}

\begin{exercise}[$\star\star$]\label{exo:pl:observability-and-incident-response:4}
Pourquoi les seuils de trente et soixante secondes ne voient-ils jamais les arrêts intermittents~?
\end{exercise}

\begin{exercise}[$\star\star$]\label{exo:pl:observability-and-incident-response:5}
Qu'apporte la fenêtre courte d'une règle de vitesse de consommation à la fenêtre longue~?
\end{exercise}

\begin{exercise}[$\star\star$]\label{exo:pl:observability-and-incident-response:6}
Quel est le \omterm{def:pl:observability-and-incident-response:ttd}{délai de rétablissement} de l'arrêt rejoué, et pourquoi ce n'est pas un temps moyen de réparation (MTTR)~?
\end{exercise}

\begin{exercise}[$\star\star\star$]\label{exo:pl:observability-and-incident-response:7}
\emph{Programmation.} Rendre les pauses deux fois plus fréquentes (\texttt{pauses\_per\_hour=1.6}) et recompter les fausses alertes du seuil de cinq secondes et
des règles de vitesse de consommation sur trente jours.
\end{exercise}

\begin{exercise}[$\star\star\star$]\label{exo:pl:observability-and-incident-response:8}
\emph{Trouver la faille.} «~Notre supervision était correcte~: chaque service a signalé un état sain pendant toute la panne.~»
\end{exercise}

\section{Problème~: Trois heures de silence}

\begin{problem}[{Problème du week-end --- trois heures de silence}]\label{pb:pl:observability-and-incident-response:1}
La chaîne du chapitre, le SLO, les jours calmes, les incidents et les règles.

\medskip\noindent\textbf{Part I --- Signals.}
\begin{enumerate}
\item Que répondent respectivement les \omterm{def:pl:observability-and-incident-response:metric}{métriques}, les journaux et les traces~?
\item Où un message tracé passe-t-il son temps un jour normal, et pendant le consommateur lent~?
\item Pourquoi le SLI est-il l'âge des données et non la santé du processus~?
\item Que permet le SLO, et que consomment les jours calmes~?
\item En quoi un SLO diffère-t-il d'un accord de niveau de service~?
\end{enumerate}

\medskip\noindent\textbf{Partie II --- Règles.}
\begin{enumerate}[resume]
\item Quels sont les trois incidents rejoués~?
\item Que détecte chaque seuil statique, et en combien de temps~?
\item Combien de fausses alertes par mois envoie chaque règle~?
\item Que détectent les règles de vitesse de consommation, et en combien de temps~?
\item Pourquoi la fausse alerte mensuelle de la règle de vitesse de consommation est-elle acceptable~?
\end{enumerate}

\medskip\noindent\textbf{Partie III --- Réponse.}
\begin{enumerate}[resume]
\item Que fait et que ne fait pas le \omterm{def:pl:observability-and-incident-response:ic}{responsable d'incident}~?
\item Qu'est-ce qui vient en premier dans la \omterm{def:pl:observability-and-incident-response:oncall}{procédure d'exploitation} pour une alerte d'âge de données dans une société de trading~?
\item Quelle est la chronologie de la panne rejouée~?
\item Quels sont son \omterm{def:pl:observability-and-incident-response:ttd}{délai de détection} et son \omterm{def:pl:observability-and-incident-response:ttd}{délai de rétablissement}~?
\item Que produit une \omterm{def:pl:observability-and-incident-response:blameless}{revue sans recherche de coupable}~?
\end{enumerate}

\medskip\noindent\textbf{Partie IV --- Le verdict.}
\begin{enumerate}[resume]
\item Énoncer le \emph{résultat nommé}~: le \omterm{def:pl:observability-and-incident-response:ttd}{délai de détection} d'un arrêt de flux et d'un consommateur lent sous un seuil statique et sous des \omterm{def:pl:observability-and-incident-response:burn}{alertes de vitesse de consommation} (burn rate), ainsi que les fausses alertes par mois que coûte chacun, sur l'incident rejoué et sur trente jours calmes.
sur l'incident rejoué et sur trente jours calmes.
\item Quelle règle déploieriez-vous, et avec quelle règle de ticket derrière~?
\item Qu'est-ce qui aurait fait virer au rouge les tableaux de bord de la société le matin de l'accroche~?
\item Quelles modifications la revue de l'arrêt rejoué exigerait-elle~?
\item En une phrase~: que doit signifier une alerte~?
\end{enumerate}
\end{problem}

\section{Questions d'entretien}

\begin{interviewq}[$\star$ \iqroles{développeur}]\label{iq:pl:observability-and-incident-response:1}
Quelle est la différence entre surveiller un processus et surveiller un service~?
\end{interviewq}

\begin{interviewq}[$\star\star$ \iqroles{développeur}]\label{iq:pl:observability-and-incident-response:2}
Qu'est-ce qu'un SLI, un SLO et un \omterm{def:pl:observability-and-incident-response:budget}{budget d'erreur}~? Donnez-en un pour une plateforme de données de marché.
\end{interviewq}

\begin{interviewq}[$\star\star$ \iqroles{développeur}]\label{iq:pl:observability-and-incident-response:3}
Votre équipe reçoit trente alertes par semaine et en ignore la plupart. Que changez-vous~?
\end{interviewq}

\begin{interviewq}[$\star\star$ \iqroles{développeur}]\label{iq:pl:observability-and-incident-response:4}
Un message met une minute pour aller du flux jusqu'au système de risque. Comment trouvez-vous où~?
\end{interviewq}

\begin{interviewq}[$\star\star\star$ \iqroles{développeur}]\label{iq:pl:observability-and-incident-response:5}
Concevez le processus d'\omterm{def:pl:observability-and-incident-response:observability}{observabilité} et de gestion des incidents pour la plateforme de données et de risque d'une société de trading.
\end{interviewq}
